\begin{abstract}
针对异构人工智能任务与能源网络耦合的跨区域调度问题，本文建立算—储—电协同模型。问题一采用固定起点季节基线与 Huber-Ridge 回归预测 GPU 需求，测试期预测总量为 $13155.739$ 等效 GPU，MAE 均值为 $25.837$ 等效 GPU；按实时优先和截止期排序并逐时择区，完成 $538$ 个测试及收尾时域任务调度，其中 $68$ 个在第2400—2405小时结清。问题二以任务迁移与开工换时为邻域实施有界局部搜索，完成 $50000$ 个任务的全时域调度；计及新能源外售收入后的净成本为 $-421711041.270$ 元，碳排放为 $188.761$ tCO2，时延为 $5.348$ ms，新能源利用率为 $67.960\%$。问题三建立含充放电、SOC 终端状态和购售电边界的能量流模型，按电价分位规则逐时递推；成本增加 $304223.929$ 元、绝对爬坡量增加 $247.775$ MW，说明该储能规则未改善经济性与波动性。问题四以成本、碳、时延、服务质量、新能源利用率和峰值净购电为指标，以无量纲评分实施有界局部搜索，生成 $5$ 个候选。推荐方案迁移 $94$ 个任务，净成本为 $-421426260.177$ 元，碳排放为 $0.238$ tCO2，时延为 $5.262$ ms，服务质量代理为 $99.944\%$，新能源利用率为 $67.973\%$，峰值净购电为 $1.006$ MW；约束验证值为 $1$。结果表明方案仅为确定性有限搜索下的情景可行候选，Pareto 性质与全局最优性均未认证。
\end{abstract}

\noindent\textbf{关键词：}算电协同；异构任务调度；Huber-Ridge 回归；储能能量流；有界局部搜索；多目标调度